\subsection{P030 --- Semantic prediction assistant approach applied to energy efficiency in Tertiary buildings}
\label{paper:P030}
\textbf{Authors:} Iker Esnaola-Gonzalez, Jesús Bermúdez, Izaskun Fernandez, Aitor Arnaiz\par
\textbf{Major Category:} BIM and Semantic Information\par
\textbf{Secondary Category:} AI and Machine Learning for Buildings, Building Energy and ZEB\par
\textbf{Keywords:} Semantic Web Technologies, Knowledge Discovery in Databases, Energy efficiency, Buildings, Ontology, Prediction\par
\paragraph{主な主張}
EEPSAは, ontology-driven semantic annotation, query/inference, outlier handling, attribute generationをKDDのselection, preprocessing, transformationへ組み込み, 実建物の室温予測支援においてbaselineよりprediction errorを低減した. \cite{EsnaolaGonzalez:2018SemanticPredictionAssistant}
\paragraph{新規性}
Semantic Webを単なるデータ統合に留めず, KDDの複数段階にdomain knowledgeを組み込み, 変数選定, データ品質処理, 特徴生成を一連のprediction workflowとして支援し, real-world use caseで評価した点に新規性がある.
\paragraph{位置付け}
semantic informationがprediction workflowへ直接作用する先行例であり, ontology-driven data selection, preprocessing, feature generation, prediction supportを検討する際の比較基準として位置付ける.
